\chapter{数据集构建与完整性审计}

\section{任务定义}

本项目输入单张胸部X光图像，输出阴性与肺炎两类概率。Kermany目录使用NORMAL/PNEUMONIA命名；RSNA的target=0更准确的医学含义是“没有肺炎目标框或未标为肺炎”，不能等同于完全健康；NIH的No Finding/Pneumonia来自报告文本挖掘。为保持代码接口统一，目录仍使用NORMAL和PNEUMONIA，但正文按具体来源解释标签。

\section{Kermany儿童胸片}

Kermany数据由公开数据页和Cell论文描述\cite{kermanydata2018,kermany2018}。本地逐文件检查得到5856张可读图像。官方目录train为5216张、val为16张、test为624张。训练集肺炎比例约74.3\%，test约62.5\%。官方val每类只有8张，不能稳定承担checkpoint、阈值或温度选择。

\begin{table}[H]
\centering
\caption{Kermany官方目录图像统计}
\begin{tabular}{lrrr}
\toprule 集合 & 阴性 & 肺炎 & 合计\\\midrule
train & 1341 & 3875 & 5216\\
val & 8 & 8 & 16\\
test & 234 & 390 & 624\\\midrule
合计 & 1583 & 4273 & 5856\\\bottomrule
\end{tabular}
\end{table}

\section{RSNA外部数据}

RSNA挑战数据建立在NIH胸片基础上，并由专家补充可能肺炎标注\cite{shih2019rsna}。完整挑战集规模远大于本项目实际使用部分。本项目从可访问镜像中得到1707张图像，按患者ID固定为train 1052、val 213、test 442。test包含阴性215张、肺炎227张。标签CSV中同一患者可能有多个框，数据准备脚本先聚合patientId，只要任一记录Target=1便记为肺炎，box count另行保存。

\begin{table}[H]
\centering
\caption{RSNA镜像可用子集统计}
\begin{tabular}{lrrr}
\toprule 集合 & 阴性/非肺炎 & 肺炎 & 合计\\\midrule
train & 560 & 492 & 1052\\
validation & 107 & 106 & 213\\
test & 215 & 227 & 442\\\bottomrule
\end{tabular}
\end{table}

该子集达到本项目预设的外部test每类至少200张门槛，但它不是完整挑战数据，也不是重新随机抽取的全国代表样本。RSNA train/val已参与部分适配方案开发，因此RSNA test在旧项目中存在多轮观察风险。新版把其定位为主要外部压力测试，并保留这一限制。

\section{NIH弱标签子集}

NIH ChestX-ray数据通过报告文本挖掘产生多疾病标签\cite{wang2017chestxray8}。项目从一个可用压缩包中选择65张No Finding和65张标签含Pneumonia的图像，形成92/12/26张train/val/test。test每类只有13张。其作用是检查结论在另一弱标签来源上是否出现方向性异常，不用于模型选择，也不用于宣布外部验证成功。

\section{审计方法}

完整性审计由\resultfile{scripts/audit\_integrity.py}完成，包含四层检查：
\begin{enumerate}
  \item 目录层：统计每个split和类别的图像数量；
  \item 患者层：Kermany按文件名提取肺炎和阴性subject编号；RSNA使用patientId；
  \item 文件层：计算全部图像SHA256，检查跨split完全重复与标签冲突；
  \item 结果层：读取评价JSON指向的预测CSV，重新计算accuracy、precision、recall、specificity、F1、AUC、Brier和混淆矩阵。
\end{enumerate}

审计是只读的，不自动移动原始数据，也不根据模型结果改变split。机器结果保存在\resultfile{results/integrity\_audit.json}。

\section{审计结果}

40份含预测CSV的评价JSON没有发现数值不一致；较早JSON缺少部分后来新增字段，属于schema演进。三个数据集均没有字节级相同图像跨split，也没有相同图像标签冲突。最严重的问题出现在Kermany患者层：肺炎类有170个subject ID同时出现在官方train和test。比如同一\texttt{person1}的细菌性图像在train，而病毒性图像在test。两者二分类标签相同，但患者解剖和采集特征并不独立。

\begin{table}[H]
\centering
\caption{完整性审计摘要}
\begin{tabularx}{\textwidth}{lrrrX}
\toprule 数据集 & 图像数 & 跨split患者ID & 精确重复组 & 结论\\\midrule
Kermany官方目录 & 5856 & 170 & 0 & 患者独立门禁失败\\
RSNA处理后子集 & 1707 & 0 & 0 & 当前划分门禁通过\\
NIH弱标签子集 & 130 & 0 & 0 & 完整性通过，样本量不足\\\bottomrule
\end{tabularx}
\end{table}

患者交叉不等于旧准确率必然被高估。旧官方test的人群和子类型组成可能比重新随机分组的test更难。后续结果确实显示，患者级随机划分的内部准确率反而更高。因此本文只作风险判断：官方结果不能作为严格患者独立估计；不把新旧分数差异全部归因于泄漏。

\section{患者级重划分}

项目汇总官方train/val/test全部图像，以标签和subject ID为组，按70\%/10\%/20\%分层分配，seed固定为42。同一subject的全部图像只进入一个集合。输出采用hardlink避免复制图像字节，逐图manifest记录新路径、标签、subject、原split、原路径和SHA256。

\begin{table}[H]
\centering
\caption{Kermany患者级新划分}
\begin{tabular}{lrrrrr}
\toprule & \multicolumn{3}{c}{图像数} & \multicolumn{2}{c}{患者标识数}\\
集合 & 阴性 & 肺炎 & 合计 & 阴性 & 肺炎\\\midrule
train & 1113 & 2994 & 4107 & 1011 & 1172\\
validation & 156 & 423 & 579 & 144 & 167\\
locked test & 314 & 856 & 1170 & 289 & 335\\\bottomrule
\end{tabular}
\end{table}

manifest SHA256以\texttt{ef73fd9733898e62}开头，完整值保存在划分摘要。新数据通过患者、哈希和标签三层审计。所有新版内部主结果基于此划分；旧官方结果只用于复现与敏感性讨论。

\section{数据集差异与混杂}

Kermany主要为儿童胸片，RSNA来自更复杂人群和不同挑战标注流程；Kermany本地文件为JPEG且保留原始纵横比，RSNA处理图为PNG并在预处理时统一尺寸；NIH又使用弱标签。模型面对的变化同时包括人群、疾病定义、图像格式和灰度风格，无法用一个变量解释。

\begin{table}[H]
\centering
\caption{三个数据来源的关键差异}
\begin{tabularx}{\textwidth}{lXXX}
\toprule 项目 & Kermany & RSNA子集 & NIH子集\\\midrule
主要人群 & 儿童 & 更复杂，含成人 & NIH临床胸片\\
标签来源 & 目录标签 & 专家挑战标注及框 & 报告文本挖掘\\
项目用途 & 开发与内部评价 & 外部压力测试/适配 & 弱标签敏感性分析\\
图像格式 & JPEG & DICOM镜像转换PNG & PNG\\
主要风险 & 原官方患者交叉 & 镜像非完整数据 & 样本极小、标签噪声\\\bottomrule
\end{tabularx}
\end{table}

\warningbox{数据集之间的性能差异同时反映模型泛化、人群差异、标签差异和图像处理差异。本文不把外部下降解释为单一算法失败，也不把RSNA阴性直接称为健康正常。}

\section{患者标识解析的误差边界}

Kermany没有随当前镜像提供独立患者表，患者标识由文件名推断。肺炎文件形如\texttt{person100\_bacteria\_475.jpeg}或\texttt{person100\_virus\_184.jpeg}，解析器取\texttt{person100}。阴性文件形如\texttt{IM-0115-0001.jpeg}或\texttt{NORMAL2-IM-1427-0001.jpeg}，解析器取检查编号前缀。

这种规则比图像级随机划分更严格，但不是完美患者真值。后缀可能表示同一检查的多张图，也可能表示不同检查；person编号跨细菌/病毒子目录时是否始终指同一真实患者，需要原始数据说明或作者元数据确认。由于重复编号至少表示高度相关的命名组，把它们放入同一split是保守选择。

新划分分别在两个二分类标签内对subject分层。这可以控制患者数比例，却不能保证图像数精确为70/10/20，因为每个患者图像数不同。肺炎患者多图更常见，最终图像比例会有轻微偏差。表中同时报告图像数与subject数，防止把图像比例误写为患者比例。

\section{文件哈希与感知重复}

SHA256只能发现字节完全相同文件。相同胸片若被重新压缩、缩放或改写元数据，SHA256会不同。项目又检查了文件名subject交叉，但尚未运行大规模感知哈希或图像相似度聚类。因此“精确重复为0”不能扩展成“没有近重复图像”。

未来可对灰度缩略图计算pHash，并对高相似候选人工复核。阈值过松会把同一患者的真实随访误判为重复，阈值过严又漏掉重新压缩图。感知重复审计应作为辅助，不替代患者元数据。

Hardlink构建的新数据目录与原图共享inode，不增加图像字节副本。它不会让训练集和test共享路径；每个源文件只链接到一个新split。若后续修改原文件，hardlink内容会一起变化，因此原始数据应只读，manifest哈希用于发现改变。

\section{标签转换的具体风险}

Kermany的PNEUMONIA包含细菌和病毒文件名子型，但主任务合并二者。若细菌/病毒在官方train/test比例不同，模型难度会变化。患者级随机划分按二分类标签分层，没有额外按子型分层，可能产生子型比例波动。新版test较容易可能部分来自这一点，后续应报告子型分布。

RSNA同一patientId可能有多个box记录。脚本先聚合，只要任一Target为1便为阳性，避免一张阳性图因重复行产生多个训练样本。Target=0包括``No Lung Opacity / Not Normal''等复杂情况时，统一目录名NORMAL可能造成语义简化。正文因此使用“阴性/非肺炎”。

NIH多标签图像只要Finding Labels包含Pneumonia就进入阳性，即使同时存在Effusion等标签；阴性只取精确No Finding。两类并不构成纯粹的“肺炎唯一病变 vs 健康”，且报告挖掘标签有假阴性和假阳性。它只适合作为压力测试。

\section{数据门禁的可维护性}

完整性审计被写成独立脚本而非一次性笔记。新增数据或重新划分后可以重复执行；评价结果变化时也能检查JSON与CSV是否一致。审计返回非零退出码表示正式门禁失败，但原Kermany目录应持续保留FAIL状态，不能通过忽略患者ID来“修绿”。

数据门禁还需要版本治理。manifest、summary和审计JSON应进入版本控制，小型预测文件可作为核心证据；大图像和checkpoint由哈希与release管理。报告中引用绝对本机路径会降低可移植性，最终公开结果应同时提供仓库相对路径。
